\chapter{Background and Related Work}

\section{Monte Carlo Simulation}

Monte Carlo (MC) simulation is a widely used technique for estimating expectations of stochastic systems through repeated random sampling~\cite{glasserman2004}. It is particularly valuable in settings where analytical solutions are unavailable or intractable, such as the pricing of financial derivatives under complex dynamics. In such applications, MC methods approximate quantities of interest by simulating a large number of independent sample paths and averaging the resulting payoffs.

A key property of MC simulation is that its statistical error decreases at a rate proportional to $O(N^{-1/2})$, where $N$ is the number of samples. As a result, achieving high accuracy typically requires a large number of simulations, making computational efficiency a primary concern. At the same time, MC workloads are often described as ``embarrassingly parallel'', since individual simulation paths can be evaluated independently. This makes them well-suited to modern parallel hardware, including multi-core CPUs and GPUs.

In practice, Monte Carlo \emph{engines} extend basic simulations by providing configurable, reusable implementations that support parameter variation, reproducibility through controlled random number generation, and repeated execution for statistical aggregation. Benchmarking such engines therefore requires not only measuring raw execution time, but also understanding how performance scales with workload size, parallelism, and implementation choices.

A useful distinction for this project is between a Monte Carlo engine and a
benchmarking framework. An engine implements the numerical workload itself: for
example, a European option-pricing engine may accept a strike, volatility,
interest rate, maturity, path count, and random seed, then return a price and
possibly sensitivities. A framework sits around such engines. It is responsible
for constructing repeated runs, changing parameters systematically, measuring
runtime and memory behaviour, recording hardware metadata, and producing
analysis artefacts. This distinction is important because many performance
claims about Monte Carlo code are not caused by the payoff formula alone, but
by the surrounding machinery: random-number generation, warm-up policy, thread
configuration, storage, and postprocessing.

\section{Automatic Differentiation}

Automatic differentiation (AD) is a technique for computing derivatives of numerical programs by systematically applying the chain rule to the sequence of elementary operations that define the computation~\cite{griewank2008,baydin2018}. Unlike finite difference methods, AD produces derivatives accurate up to floating-point precision while avoiding numerical instability.

Two primary modes of AD are commonly used. Forward-mode AD propagates derivatives alongside values during the forward execution of a program and is typically efficient when the number of input parameters is small relative to the number of outputs. In contrast, reverse-mode AD records intermediate computations during a forward pass and then propagates sensitivities backwards, making it more efficient when a small number of outputs depend on many inputs.

Modern AD systems, such as those integrated into scientific computing and machine learning frameworks, enable gradients to be computed directly from program implementations. JAX is one example of this style of system, combining NumPy-like array programming with composable transformations such as JIT compilation and automatic differentiation~\cite{jax2018}. This is particularly useful in simulation-based contexts, where sensitivities (e.g., financial ``Greeks'') can be obtained without manual derivation.

However, the choice between forward- and reverse-mode AD introduces important trade-offs in computational cost and memory usage. Forward-mode typically incurs moderate overhead with minimal additional memory requirements, while reverse-mode can significantly increase memory usage due to the need to store intermediate states. These trade-offs become especially important in large-scale simulations.

Several AD systems are relevant when considering how MC+AD support might be
extended beyond the JAX implementation evaluated in this project. Enzyme
differentiates LLVM intermediate representation, making it attractive for C and
C++ style numerical kernels~\cite{moses2021enzyme}. Zygote provides
source-to-source AD in the Julia ecosystem~\cite{innes2019zygote}, while
DiffTaichi demonstrates differentiable simulation on accelerator-oriented
programming abstractions~\cite{hu2020difftaichi}. These systems illustrate
that AD is not a single implementation technique. The design choice for this
project was therefore to treat AD mode as part of the benchmark configuration,
rather than assuming that all engines expose identical derivative capability.


\section{Monte Carlo Greeks in Finance}

In quantitative finance, sensitivities such as Delta, Vega and Rho are often as
important as the option price itself. Standard Monte Carlo Greek estimators
include finite differences, pathwise derivative methods, and likelihood-ratio
or score-function methods~\cite{glasserman2004,giles2006smoking}. Finite
differences are simple but introduce step-size and noise trade-offs; common
random numbers reduce variance but do not remove the bias-variance choice.
Pathwise and likelihood-ratio estimators avoid some of these problems under
appropriate smoothness or density assumptions.

Adjoint algorithmic differentiation extends this idea computationally: once a
pricing computation is written as a program, reverse accumulation can produce
many sensitivities at a cost that is often a small multiple of the primal
valuation~\cite{capriotti2011fast}. This project does not propose a new Greek
estimator. Existing work studies Monte Carlo Greeks and AD/AAD methods, while
this project studies the systems-level benchmarking, reproducibility and
cost-performance behaviour of implemented MC+AD workloads across backends and
cloud CPU configurations.

\section{Monte Carlo + AD Challenges}

Combining Monte Carlo simulation with automatic differentiation enables flexible and general sensitivity analysis, but introduces several non-trivial computational challenges.

First, the computational cost of AD compounds with the cost of simulation. In MC workloads, each sample path may consist of many time steps and numerical operations. Applying AD to such computations increases the amount of work performed per path, with the overhead depending on the differentiation mode and implementation. As a result, throughput (e.g., paths per second) can be significantly reduced compared to non-differentiated simulations.

Second, memory usage becomes a critical concern, particularly for reverse-mode AD. Reverse-mode requires storing intermediate values during the forward execution in order to compute gradients during the backward pass. In long-horizon simulations with many time steps, this can lead to substantial memory growth, potentially exceeding cache capacity and degrading performance due to increased memory traffic. In extreme cases, memory limitations may prevent scaling to larger workloads altogether.

Third, the interaction between AD and parallelism is complex. While Monte Carlo simulation is naturally parallel across independent paths, AD introduces dependencies within each path that may limit opportunities for parallel execution. Furthermore, reverse-mode AD can introduce additional synchronisation points and data movement, particularly in GPU or distributed settings, reducing the effectiveness of parallelisation strategies.

Finally, the combination of stochastic computation and differentiation introduces challenges for reproducibility and numerical stability. Ensuring that differentiated simulations produce consistent and meaningful results requires careful control of random number generation, seeding, and execution order.

These challenges motivate the need for systematic evaluation of MC+AD workloads, where both computational performance and memory behaviour are measured under controlled and reproducible conditions.

\section{Existing Frameworks}

Several frameworks exist for Monte Carlo simulation and high-performance
numerical computing, though few target MC+AD benchmarking across heterogeneous
backends. Carlo.jl is a relevant example because it separates the Monte Carlo
algorithm from scheduling, checkpointing, storage and postprocessing, and it
emphasises reproducible execution~\cite{weber2025carlojl}. Those ideas
influenced this project: workload definitions, engine execution, measurement
policy and storage are separated so that benchmark rows remain auditable when
engines are added or replaced.

However, Carlo.jl is intentionally scoped to a single language ecosystem and
does not directly address cross-language benchmarking, AD support across
backends, or cloud cost-performance. More broadly, scientific and machine
learning benchmarks often focus on linear algebra kernels or neural network
training rather than the interaction between stochastic simulation,
differentiation and parallel execution. This motivates a framework that treats
Monte Carlo simulation, AD-capable execution, reproducibility metadata and cloud
cost as first-class benchmarking concerns.

\section{Hardware Considerations}

MC and MC+AD performance depends on arithmetic throughput, memory bandwidth,
cache behaviour, vectorisation, multi-threading and runtime overheads. GPUs can
provide high throughput for large data-parallel simulations, but kernel launch,
memory transfer, control-flow divergence and AD memory requirements complicate
the comparison. In cloud environments, raw speed must also be weighed against
instance price and utilisation. The evaluation therefore reports both runtime
and cost-performance rather than assuming that the fastest machine is the best
deployment choice.

\section{Benchmarking Methodology and Reproducibility}

Benchmarking numerical workloads is difficult because small methodological
differences can change the result. Warm-up effects, JIT compilation, thread
pool initialisation, random-number generation, cache state, and background
system load can all distort a single timing. This is particularly important
for MC+AD workloads, where the cost per path may vary with the number of time
steps, the differentiation mode, and the memory layout used by the engine.

For low-level CPU studies, benchmark suites such as PolyBench/C are useful
because they isolate common loop nests and memory-access patterns under
controlled conditions~\cite{pouchet2012polybench}. MC+AD workloads are higher
level than these kernels, but the same lesson applies: the benchmark must make
the execution conditions explicit. This project therefore uses warm-up runs,
repeated timed measurements, deterministic seeds where supported, and database
records that capture both workload parameters and environment metadata.

Hardware-counter and profiler tools provide another perspective on
performance. LIKWID exposes CPU topology and performance counters in a
lightweight workflow~\cite{treibig2010likwid}, while Intel VTune is widely used
for deeper CPU performance analysis~\cite{intelvtune}. For XLA-backed systems,
OpenXLA's profiling documentation illustrates how modern compiler runtimes
surface kernel-level and memory behaviour~\cite{xprof2025}. These tools were
not used as primary evidence in the final results, but they inform the
framework design: runtime alone is not always enough to explain why an engine
scales or fails to scale, so metadata capture and profiler integration are
natural extensions of the same benchmark lifecycle.

\section{Compiler Techniques from Machine Learning Systems}

Several compiler techniques associated with machine-learning systems are
relevant to scientific simulation. Operator fusion and graph compilation can
reduce memory traffic and dispatch overhead, which matters for Monte Carlo
workloads that repeatedly apply similar update kernels. MLIR provides one
example of an extensible compiler infrastructure for lowering computations to
different hardware targets~\cite{lattner2020mlir}. TVM demonstrates automated
code generation and optimisation for tensor programs~\cite{chen2018tvm}, and
Ansor extends this direction with hardware-aware search~\cite{zheng2020ansor}.

The relevance to this project is not that MC+AD workloads are identical to
neural-network training, but that both domains benefit from expressing
computation in a form that compilers can optimise. In the implementation
chapter, this appears concretely in the use of JAX/XLA and \texttt{jax.lax.scan}
for the local-volatility loop. Expressing the loop as a staged computation
allows XLA to compile it as a structured loop rather than repeatedly dispatching
Python operations. The evaluation therefore treats JAX/XLA as an AD-capable
compiled runtime whose behaviour should be measured against NumPy and native
compiled engines, rather than assumed to be faster or slower by design.

Mixed precision is another technique that has been successful in machine
learning workloads~\cite{micikevicius2017mixed}. For Monte Carlo simulation,
the question is more delicate because reducing precision can affect both
statistical estimates and pathwise derivatives. This project records precision
and backend metadata where available, but the final evaluation focuses on
double-precision CPU-oriented runs. The broader point is that benchmark
configuration should expose such choices explicitly instead of hiding them as
engine-specific optimisations.

\section{Benchmark Reporting Standards}

Large-scale benchmarks such as DAWNBench and MLPerf emphasise transparent
reporting of workload, hardware, software versions and measurement protocol
\cite{coleman2017dawnbench,coleman2019dawnbench,mattson2020mlperf}. This
report applies the same discipline at project scale: tables identify workload,
engine, AD mode, path count and instance type, while the experiments chapter
states warm-up, repetition, validation and exclusion policies.

\section{Budget-Aware Benchmark Selection}

Full-grid benchmarking is attractive because it gives a direct comparison of
all candidate workloads, engines, AD modes, and hardware choices. However, it
becomes expensive when each candidate must be evaluated at production-scale path
counts on cloud instances. A practical benchmarking system therefore needs not
only accurate measurement, but also a way to decide which configurations deserve
the most expensive measurements.

This motivates the use of cheap probe runs. A probe run evaluates a candidate
configuration at a smaller path count or reduced budget and uses that result as
evidence about whether the configuration is likely to remain competitive at full
scale. The central risk is selection error: a configuration that is slow at low
budget may become competitive at high budget, or a noisy probe may give an
incorrect ranking. Any profiler must therefore be evaluated against a full-grid
baseline rather than assumed to be correct.

Successive halving is a natural fit for this setting. It begins by evaluating
all candidates at a low budget, promotes a smaller active set to a larger budget,
and repeats until only the strongest candidates remain for full evaluation. In
the context of MC+AD benchmarking, the budget can be interpreted as the number
of simulated paths. This connects the statistical structure of Monte Carlo
experiments with the systems problem of cloud cost control: most configurations
receive only cheap measurements, while promising configurations receive enough
budget to support reliable final decisions.

\section{Summary of Gaps}

The literature motivates three gaps addressed by this project. First, existing
Monte Carlo frameworks rarely combine reproducible MC execution, AD-capable
workloads, cross-backend comparison and cloud metadata in one benchmark loop.
Second, AD performance in stochastic simulation is less standardised than in
machine-learning benchmarks, particularly for memory pressure and parallel
execution. Third, practical deployment requires cost-performance and profiler
selection evidence, not only raw runtime. This project addresses these gaps by
building a modular benchmarking framework and using it to evaluate implemented
MC+AD workloads across software stacks and cloud CPU configurations.
